% GENERATED FILE. Edit canonical lessons, then run npm run generate:books.
% canonical-source-hash: fnv1a64:8fa9d5324a5aea62
% canonical-lessons: TE-C103-gudi, TE-C103-badi, TE-C103-grandhalayam, TE-C103-byanku, TE-C103-santa

\chapter{Temple, School, Library, Bank, Market}
\label{ch:te-temple-school-library-bank-market}
\hlchaptermodality{103}

\begin{chapteropening}
\textbf{By the end of this chapter:} \emph{I can say a temple, a school, a library, a bank and a market.}

Complete the last lesson of chapter 103: I can say a temple, a school, a library, a bank and a market.
\end{chapteropening}

\section[\texorpdfstring{\te{గుడి} (guḍi)}{గుడి (guḍi)}]{\te{గుడి} (guḍi) — a temple}

\label{lesson:TE-C103-gudi}



\begin{quote}
Before the new one: say the Telugu for a bridge, then the Telugu for a park.
\end{quote}

\subsection*{You'll want to know: \te{గుడి}}
\textbf{\te{గుడి}} — \emph{guḍi} — \textquotedblleft{}a temple\textquotedblright{}.

\begin{grammarlens}[title={What you've built}]
One more word for this chapter.
\end{grammarlens}

\subsection*{Your turn}
\begin{itemize}
\raggedright
  \item \emph{Say it:} \emph{guḍi}
  \item \emph{Say it:} \emph{guḍi}, once more
  \item \emph{Say it:} say \emph{guḍi}
  \item \emph{Recall:} say the Telugu for a bridge, then the Telugu for a park, then say \emph{guḍi} again
\end{itemize}

\begin{tcolorbox}[breakable,colback=teal!4,colframe=teal!35!black,title={Before you move on}]
What does \te{గుడి} mean? (\textquotedblleft{}A temple\textquotedblright{}.) Say it once more.
\end{tcolorbox}
\section[\texorpdfstring{\te{బడి} (baḍi)}{బడి (baḍi)}]{\te{బడి} (baḍi) — a school}

\label{lesson:TE-C103-badi}



\begin{quote}
Before the new one: say the Telugu for a park, then the Telugu for a temple.
\end{quote}

\subsection*{You'll want to know: \te{బడి}}
\textbf{\te{బడి}} — \emph{baḍi} — \textquotedblleft{}a school\textquotedblright{}.

\begin{grammarlens}[title={What you've built}]
One more word for this chapter.
\end{grammarlens}

\subsection*{Your turn}
\begin{itemize}
\raggedright
  \item \emph{Say it:} \emph{baḍi}
  \item \emph{Say it:} \emph{baḍi}, once more
  \item \emph{Say it:} say \emph{baḍi}
  \item \emph{Recall:} say the Telugu for a park, then the Telugu for a temple, then say \emph{baḍi} again
\end{itemize}

\begin{tcolorbox}[breakable,colback=teal!4,colframe=teal!35!black,title={Before you move on}]
What does \te{బడి} mean? (\textquotedblleft{}A school\textquotedblright{}.) Say it once more.
\end{tcolorbox}
\section[\texorpdfstring{\te{గ్రంథాలయం} (granthālayaṁ)}{గ్రంథాలయం (granthālayaṁ)}]{\te{గ్రంథాలయం} (granthālayaṁ) — a library}

\label{lesson:TE-C103-grandhalayam}



\begin{quote}
Before the new one: say the Telugu for a temple, then the Telugu for a school.
\end{quote}

\subsection*{You'll want to know: \te{గ్రంథాలయం}}
\textbf{\te{గ్రంథాలయం}} — \emph{granthālayaṁ} — \textquotedblleft{}a library\textquotedblright{}.

\begin{grammarlens}[title={What you've built}]
One more word for this chapter.
\end{grammarlens}

\subsection*{Your turn}
\begin{itemize}
\raggedright
  \item \emph{Say it:} \emph{granthālayaṁ}
  \item \emph{Say it:} \emph{granthālayaṁ}, once more
  \item \emph{Say it:} say \emph{granthālayaṁ}
  \item \emph{Recall:} say the Telugu for a temple, then the Telugu for a school, then say \emph{granthālayaṁ} again
\end{itemize}

\begin{tcolorbox}[breakable,colback=teal!4,colframe=teal!35!black,title={Before you move on}]
What does \te{గ్రంథాలయం} mean? (\textquotedblleft{}A library\textquotedblright{}.) Say it once more.
\end{tcolorbox}
\section[\texorpdfstring{\te{బ్యాంకు} (byāṅku)}{బ్యాంకు (byāṅku)}]{\te{బ్యాంకు} (byāṅku) — a bank}

\label{lesson:TE-C103-byanku}



\begin{quote}
Before the new one: say the Telugu for a school, then the Telugu for a library.
\end{quote}

\subsection*{You'll want to know: \te{బ్యాంకు}}
\textbf{\te{బ్యాంకు}} — \emph{byāṅku} — \textquotedblleft{}a bank\textquotedblright{}.

\begin{grammarlens}[title={What you've built}]
One more word for this chapter.
\end{grammarlens}

\subsection*{Your turn}
\begin{itemize}
\raggedright
  \item \emph{Say it:} \emph{byāṅku}
  \item \emph{Say it:} \emph{byāṅku}, once more
  \item \emph{Say it:} say \emph{byāṅku}
  \item \emph{Recall:} say the Telugu for a school, then the Telugu for a library, then say \emph{byāṅku} again
\end{itemize}

\begin{tcolorbox}[breakable,colback=teal!4,colframe=teal!35!black,title={Before you move on}]
What does \te{బ్యాంకు} mean? (\textquotedblleft{}A bank\textquotedblright{}.) Say it once more.
\end{tcolorbox}
\section[\texorpdfstring{\te{సంత} (santa)}{సంత (santa)}]{\te{సంత} (santa) — a market}

\label{lesson:TE-C103-santa}



\begin{quote}
Before the new one: say the Telugu for a library, then the Telugu for a bank.
\end{quote}

\subsection*{You'll want to know: \te{సంత}}
\textbf{\te{సంత}} — \emph{santa} — \textquotedblleft{}a market\textquotedblright{}.

\begin{grammarlens}[title={What you've built}]
One more word for this chapter.
\end{grammarlens}

\subsection*{Your turn}
\begin{itemize}
\raggedright
  \item \emph{Say it:} \emph{santa}
  \item \emph{Say it:} \emph{santa}, once more
  \item \emph{Say it:} say \emph{santa}
  \item \emph{Recall:} say the Telugu for a library, then the Telugu for a bank, then say \emph{santa} again
\end{itemize}

\begin{tcolorbox}[breakable,colback=teal!4,colframe=teal!35!black,title={Before you move on}]
What does \te{సంత} mean? (\textquotedblleft{}A market\textquotedblright{}.) Say it once more.
\end{tcolorbox}
